\section{Introduction}
\label{sec:intro}

Uncrewed aerial imagery is now a routine instrument for wide-area search
and rescue, and automated person detection is what makes reviewing that
coverage tractable. Speed dominates outcomes: incident data show most
missing subjects are located within the first day, with recovery odds
declining as search time lengthens~\cite{koester2008lost}. A detector
that misses a person in frame is not a benign error.

The two standard benchmarks for this task, HERIDAL~\cite{bozic2019heridal}
and SARD~\cite{sambolek2021sard}, both report strong accuracy, and both
were collected in calm, clear, well-lit conditions. Real missions are
launched because something has gone wrong: after a flood, during a
wildfire, in a storm, or at night. No existing benchmark systematically
measures how these detectors behave under the atmospheric conditions
that trigger their deployment, so the gap between benchmarked accuracy
and operational reliability is, at present, unmeasured.

This gap is unusually tractable in Earth observation, because the
relevant conditions are physically characterisable. Atmospheric
scattering, sediment-laden water, rain, and photon-limited imaging all
have established optical models in the remote-sensing and vision
literature~\cite{koschmieder1924theorie,narasimhan2002vision,garg2007rain,chen2018dark}.
Rather than argue that the evaluation is missing, we build it from those
models.

AegisBench applies nine physically modelled corruptions across four
disaster families to the HERIDAL and SARD test sets and evaluates three
architecturally distinct detectors over all resulting conditions. Three
choices make the result an evaluation protocol rather than a
demonstration. Each corruption is derived from a documented optical
model of a named disaster condition, so a result maps onto an operating
condition a responder can recognise. Each severity is calibrated against
a declared, machine-verified image statistic rather than chosen by eye,
so severity claims are falsifiable. Every detector's confidence
threshold is selected once on clean validation data and frozen across
all corrupted conditions, matching how a system is fielded rather than a
best case unavailable mid-mission.

What the benchmark then measures is not a uniform loss of accuracy. It
is a set of failures that differ by condition, by capture regime and by
architecture, with one of them sitting at the absolute floor. We report
the measurement, establish that the floor is not an artefact of
thresholds, seeds or the corruption model, and then show what it takes
to move it.

\paragraph{Contributions.}
\begin{itemize}\itemsep1pt \parskip0pt \topsep2pt
\item A physically grounded corruption benchmark for aerial
search-and-rescue person detection: nine corruptions, four disaster
families, three calibrated severities, bit-reproducible from the source
datasets without redistributing imagery
(Section~\ref{sec:benchmark}).
\item A 168-condition evaluation of three architecturally distinct
detectors under one shared evaluator at frozen operating points, with
bootstrap intervals on every condition, showing that degradation is
uneven, that it does not transfer between capture regimes, and that
low light drives all three architectures to exactly zero recall
(Section~\ref{sec:results}).
\item Three independent checks that the zero is real rather than an
artefact: a threshold-independent mAP analysis, a three-seed repeat, and
an evaluation on real night drone imagery at the frozen operating point
(Sections~\ref{sec:results} and~\ref{sec:mitigation}).
\item A localisation analysis showing the failure is a loss of
detections rather than of box placement, which identifies the failure as
silent and locates where a mitigation must act
(Section~\ref{sec:results}).
\item A mitigation study that lifts severity-3 low-light recall off the
floor by training only on severities 1 and 2, and a measurement of what
that costs on the other eight corruptions, which turns out to be
architecture-dependent rather than a single robustness tax
(Section~\ref{sec:mitigation}).
\end{itemize}
